\documentclass[tikz,border=18pt]{standalone}
\usepackage{luatexja}
\usepackage{tikz}
\usetikzlibrary{shapes.geometric,arrows.meta,positioning,calc}

\definecolor{hwcopper}{RGB}{190, 80, 20}
\definecolor{alertred}{RGB}{200, 30, 40}
\definecolor{shieldgreen}{RGB}{20, 120, 60}

\begin{document}
\begin{tikzpicture}[>=Stealth, font=\small]

  % Title
  \node[draw=hwcopper, fill=orange!8, rounded corners=8pt, line width=1.5pt, minimum width=21.0cm, minimum height=1.2cm, align=center] at (0, 7.5) {
    \large\bfseries ハードウェア・セキュリティ：物理デバッグ端子（UART/JTAG）とサイドチャネル攻撃（DPA）
  };

  % Left Box: Hardware Interface Hacking
  \node[draw=hwcopper, fill=orange!5, rounded corners=8pt, line width=1.3pt, minimum width=9.8cm, minimum height=6.2cm, align=center] at (-5.3, 3.2) {};
  \node[font=\bfseries\large, text=hwcopper!90!black] at (-5.3, 5.8) {🔌 物理インターフェースの攻撃（基板ハッキング）};

  \node[draw=hwcopper, fill=white, rounded corners=4pt, thick, minimum width=9.0cm, minimum height=1.6cm, align=left] at (-5.3, 4.4) {
    \textbf{1. デバッグ端子の放置（UART / JTAGポート）}\\
    \footnotesize ・基板上のテストピン（GND, TX, RX, VCC）にUSBシリアルを直結\\
    \footnotesize ・管理者Rootシェルがパスワードなしで立ち上がる脆弱性
  };

  \node[draw=hwcopper, fill=white, rounded corners=4pt, thick, minimum width=9.0cm, minimum height=1.6cm, align=left] at (-5.3, 2.2) {
    \textbf{2. フラッシュメモリ（SPI Flash）の直接ダンプ}\\
    \footnotesize ・基板上のROMチップをクリップで挟み、ファームウェアバイナリを抽出\\
    \footnotesize ・ハードコードされた暗号鍵やパスワードハッシュを直接復元
  };

  % Right Box: Side-Channel Attacks
  \node[draw=alertred, fill=red!4, rounded corners=8pt, line width=1.3pt, minimum width=9.8cm, minimum height=6.2cm, align=center] at (5.3, 3.2) {};
  \node[font=\bfseries\large, text=alertred] at (5.3, 5.8) {⚡ サイドチャネル攻撃（電力解析・電磁波漏洩）};

  \node[draw=alertred, fill=white, rounded corners=4pt, thick, minimum width=9.0cm, minimum height=1.6cm, align=left] at (5.3, 4.4) {
    \textbf{1. 差分電力解析（DPA: Differential Power Analysis）}\\
    \footnotesize ・CPUが「0」と「1」を計算する際の微小な消費電力の揺らぎをオシロスコープで計測\\
    \footnotesize ・数学的に破れないはずのAES/RSA暗号鍵の秘密ビットを統計的に特定
  };

  \node[draw=shieldgreen, fill=green!10, rounded corners=4pt, thick, minimum width=9.0cm, minimum height=1.6cm, align=left] at (5.3, 2.2) {
    \textbf{🛡️ ハードウェア防衛策}\\
    \footnotesize ・\textbf{端子無効化・ヒューズ溶断}（製品出荷時のJTAG無効化）\\
    \footnotesize ・\textbf{セキュアブート}（公開鍵署名で改ざんファームウェアの起動拒否）\\
    \footnotesize ・\textbf{マスキング・電力ダミー消費}（サイドチャネル対策）
  };

  % Bottom note
  \node[draw=hwcopper, fill=orange!10, rounded corners=6pt, line width=1.2pt, minimum width=21.0cm, minimum height=0.9cm, align=center] at (0, -0.4) {
    \bfseries 物理的なアクセスを許せばソフトウェアの壁は突破される：ハードウェアRoot of Trust（RoT）の重要性
  };

\end{tikzpicture}
\end{document}
